% Recuperación del producto nativo y su selector, no incorporación de RH.
\subsection{Irreducibilidad multiplicativa y períodos de retorno}
\label{primos-retornos}

La distinción entre estructura, evaluación y fase aparece también en el
sector multiplicativo. El producto local de APP contiene tanto el residuo
como el cociente; su iteración permite construir formas normales y
factorizaciones antes de examinar sus períodos decimales. Este desarrollo
sitúa la expresión \emph{modo primo} en un dominio preciso y relaciona la
irreducibilidad con las operaciones ya definidas, no con una coincidencia
aislada de períodos \cite{hmtintegral,hmtholonomia}.

Sea
\[
 \mathcal W_9=\{\varnothing\}\sqcup
 \bigsqcup_{\ell\geq1}\{1,\ldots,9\}^{\ell}.
\]
Las secuencias se ordenan desde el dígito menos significativo. La evaluación
\(\operatorname{ev}_9(u)=\sum_j u_j9^j\), con
\(\operatorname{ev}_9(\varnothing)=0\), es biyectiva sobre los enteros no
negativos: para un entero positivo, la división de \(n-1\) por \(9\)
determina de manera única el primer dígito y la continuación. Esta es la
numeración nonádica biyectiva; conserva el representante \(9\) de la
frontera.

La suma \(\boxplus_\#\) se obtiene aplicando la hoja aditiva de APP a los
dígitos de cada posición y propagando sus cocientes. Si sólo hay un dígito,
éste constituye el residuo provisional; un cociente entrante no nulo se
normaliza mediante una segunda operación aditiva. La suma de los
cocientes se transporta a la posición siguiente. Para multiplicar una
secuencia por un dígito \(d\), la celda multiplicativa produce
\(u_jd=9q_j+r_j\), y el cociente entrante \(c_j\) se incorpora mediante
\[
 (w_j,c_{j+1})=
 \begin{cases}
 (r_j,q_j),&c_j=0,\\
 \bigl(R^+(r_j,c_j),q_j+q^+(r_j,c_j)\bigr),&c_j>0.
 \end{cases}
\]
Se comienza con \(c_0=0\) y se incorpora el cociente final cuando es no
nulo. Durante la suma, \(c_j\leq2\); durante el producto por un dígito,
\(c_j\leq9\). Las operaciones locales reciben así argumentos dentro
de su dominio. Denotemos por \(d\odot_\#u\) este segundo procedimiento,
con \(d\odot_\#\varnothing=\varnothing\).

Para \(v=(v_0,\ldots,v_{m-1})\), el producto completo se construye por
la recurrencia
\[
 a_m=\varnothing,\qquad
 a_j=(9\odot_\#a_{j+1})\boxplus_\#(v_j\odot_\#u),
 \qquad u\boxtimes_\#v=a_0.
\]
La regla compone las dos hojas locales. Su evaluación verifica
\begin{equation}
 \operatorname{ev}_9(u\boxtimes_\#v)
   =\operatorname{ev}_9(u)\operatorname{ev}_9(v).
 \label{eq:primos-entrelazamiento}
\end{equation}
En efecto, cada transición conserva la suma parcial más el cociente
pendiente ponderado por la potencia de \(9\). La recurrencia da entonces
\(\operatorname{ev}_9(a_j)=9\operatorname{ev}_9(a_{j+1})+
v_j\operatorname{ev}_9(u)\), cuya iteración prueba
\eqref{eq:primos-entrelazamiento}. La unicidad de la forma normal transporta
asociatividad y conmutatividad al producto construido; su unidad es \((1)\).

\begin{definition}[Coproducto multiplicativo reducido]
Sobre el espacio vectorial de soporte finito generado por las formas
normales no vacías se define
\[
 \widetilde\Delta_\# e_w
 =\sum_{\substack{u\boxtimes_\#v=w\\u,v\ne(1)}}e_u\otimes e_v.
\]
Una forma normal no unitaria es irreducible cuando su imagen por
\(\widetilde\Delta_\#\) es nula.
\end{definition}

Las fibras de factorización son finitas por
\eqref{eq:primos-entrelazamiento}. En la completación
\(\mathcal H_\#=\ell^2(\mathcal W_9\setminus\{\varnothing\})\),
las imágenes de formas distintas tienen soportes disjuntos. Si
\(d_\#(w)\) cuenta las factorizaciones ordenadas no unitarias, el dominio
de la clausura es
\[
 \operatorname{Dom}(\overline{\widetilde\Delta_\#})
 =\left\{x\in\mathcal H_\#:
       \sum_w d_\#(w)|x_w|^2<\infty\right\}.
\]
En efecto, la norma cuadrada de la imagen es esa suma ponderada. Por ello,
el operador
\[
 A_\#=(\overline{\widetilde\Delta_\#})^*
                  \overline{\widetilde\Delta_\#}
\]
es diagonal, positivo y autoadjunto: su valor propio en \(e_w\) es el
número de factorizaciones ordenadas no unitarias de \(w\). La proyección
\[
 P_{\rm irr}=(I-|e_{(1)}\rangle\langle e_{(1)}|)
                    \mathbf1_{\{0\}}(A_\#)
\]
selecciona exactamente las formas irreducibles. Mediante la evaluación
multiplicativa, éstas corresponden a los primos ordinarios. La selección
depende de la estructura de factorización, y mantiene separado el estado
unitario.

Para un valor \(n\geq2\) coprimo con \(10\), la división decimal induce la
permutación de residuos \(r\mapsto10r\pmod n\). El retorno del residuo
\(1\) tiene período
\[
 \tau_n=\operatorname{ord}_n(10).
\]
Si \(n=\prod_p p^{a_p}\), el teorema chino de los restos da
\[
 \tau_n=\operatorname{mcm}_{p^{a_p}\parallel n}
                    \operatorname{ord}_{p^{a_p}}(10).
\]
El retorno compuesto sincroniza los retornos de las potencias primas.
Para un primo \(p\notin\{2,3,5\}\), se tiene
\[
 \tau_p\mid p-1,\qquad
 M_p=\frac{10^{\tau_p}-1}{p}\in\mathbb N,
 \qquad M_p\equiv0\pmod9.
\]
La última congruencia expresa el cierre nonádico de la repetición decimal:
\(10^{\tau_p}-1\) es múltiplo de \(9\) y \(p\) es invertible módulo
\(9\). La misma congruencia vale para cualquier \(n\) coprimo con
\(90\). En particular, \(7,13,77,91,143\) tienen período decimal \(6\),
aunque sólo los dos primeros son irreducibles. El coproducto y el período
registran, por tanto, propiedades diferentes de un mismo valor construido.

Esta diferenciación precisa la relación con el refinamiento aperiódico.
La secuencia de vacancias determina los incrementos del índice de
capacidad; la factorización determina los modos multiplicativos; la
división decimal determina sus tiempos de retorno. El representante
nonádico conserva el cierre de fase, mientras que el cociente y la
factorización conservan información que ese cierre no distingue. El
enlace con las construcciones espectrales del tratado reside en esos
operadores y sus entrelazamientos. La demostración de sus resultados
analíticos posteriores mantiene su desarrollo propio; no se sustituye por
la periodicidad de una observación residual.

\subsubsection{Discrepancia de vacancias y lecturas de Dirichlet}

El vínculo admite también una expresión analítica exacta. Se conserva la
misma codificación \(z_n=z_n(0)\), para \(n\geq1\), y se consideran
dos observables posteriores: la serie sobre todos los índices y el
producto sobre los índices seleccionados por irreducibilidad. Para
\(\operatorname{Re}s>1\),
\[
 Z_{\rm vac}^{+}(s)=\sum_{n\geq1}\frac{z_n}{n^s},\qquad
 Z_{\rm vac}^{\times}(s)=\prod_p(1-p^{-s})^{-z_p}.
\]
Ambos convergen absolutamente en ese semiplano. Su factor común es la
densidad \(\delta\), pero sus discrepancias son diferentes.

\Needspace{21\baselineskip}
\begin{proposition}[Separación de la densidad de vacancias]
Se tiene
\[
 Z_{\rm vac}^{+}(s)=\delta\zeta(s)+H_{\rm vac}(s),
 \qquad H_{\rm vac}\in\mathcal O(\operatorname{Re}s>0).
\]
En \(\operatorname{Re}s>1\), con logaritmos definidos por las series
absolutamente convergentes de Euler,
\[
 Z_{\rm vac}^{\times}(s)=\zeta(s)^\delta G_{\rm vac}(s),\qquad
 \log G_{\rm vac}(s)=
 \sum_p(z_p-\delta)\log(1-p^{-s})^{-1}.
\]
\end{proposition}
\begin{proof}
Las sumas parciales centradas satisfacen
\[
 A_N:=\sum_{n=1}^N(z_n-\delta)
 =\lfloor(N+1)\delta\rfloor-N\delta,\qquad |A_N|<1.
\]
La sumación de Abel representa el término centrado por
\(s\int_1^\infty A_{\lfloor x\rfloor}x^{-s-1}\,dx\). La integral
converge localmente de manera uniforme en \(\operatorname{Re}s>0\),
lo que prueba la holomorfía. La segunda identidad resulta de separar
\(z_p=\delta+(z_p-\delta)\) en el logaritmo del producto, dentro de
su dominio de convergencia absoluta.
\end{proof}

La descomposición hace visible qué añade la lectura sobre primos: examina
la misma sucesión de vacancias sobre un subconjunto determinado por la
factorización. El control de la discrepancia sobre todos los enteros da
la primera prolongación holomorfa; el segundo observable conserva una
discrepancia propia sobre primos. Son dos aplicaciones de la misma
estructura discreta, con dominios analíticos expresos. Esta articulación
explica su pertinencia en el núcleo aritmético sin trasladar a este
artículo la demostración de los resultados terminales del tratado.
